% GENERATED FILE. Edit canonical lessons, then run npm run generate:books.
% canonical-source-hash: fnv1a64:77f4ef8a7ea00e70
% canonical-lessons: ZH-C508-jiemei, ZH-C508-xiongmei, ZH-C508-dushengzinu, ZH-C508-qifu, ZH-C508-dongshi, ZH-R508-review-504-508

\chapter{Brothers and Sisters}
\label{ch:zh-brothers-and-sisters}
\hlchaptermodality{508}

\begin{chapteropening}
\textbf{By the end of this chapter:} \emph{I can talk about sisters, a brother and sister, only children, being picked on and growing up sensible.}

Complete the last lesson of chapter 508: I can talk about sisters, a brother and sister, only children, being picked on and growing up sensible.
\end{chapteropening}

\section[jiěmèi]{jiěmèi — sisters}

\label{lesson:ZH-C508-jiemei}



\begin{quote}
Before the new one: say the Mandarin for a married couple, then the Mandarin for one's elders.
\end{quote}

\subsection*{You'll want to know: jiěmèi}
\emph{jiěmèi} — \textquotedblleft{}sisters\textquotedblright{}.

\begin{quote}
\emph{wǒmen jiěmèi sān ge wǒ zuì xiǎo} — There are three of us sisters, and I'm the youngest.
\end{quote}

Its characters come later; for now it is a word for the ear.

\begin{cousinweb}
\emph{jiěmèi} is built from \emph{jiě} 'an elder sister' and \emph{mèi} 'a younger sister', both inherited from Middle Chinese, \emph{tsjaeX} and \emph{mwojH} (Guangyun, in Baxter's transcription), which glosses the first as a word for mother among a people of the west and the second 'a younger sister'; the sense 'elder sister' is not given for the first there.
\end{cousinweb}

\begin{grammarlens}[title={What you've built}]
One more word for saying how it was.
\end{grammarlens}

\subsection*{Your turn}
\begin{itemize}
\raggedright
  \item \emph{Say it:} \emph{jiěmèi}
  \item \emph{Say it:} \emph{jiěmèi}, once more
  \item \emph{Say it:} \emph{wǒmen jiěmèi sān ge wǒ zuì xiǎo} — There are three of us sisters, and I'm the youngest.
  \item \emph{Recall:} say the Mandarin for a married couple, then the Mandarin for one's elders, then say \emph{jiěmèi} again
\end{itemize}

\begin{tcolorbox}[breakable,colback=teal!4,colframe=teal!35!black,title={Before you move on}]
What does jiěmèi mean? (\textquotedblleft{}Sisters\textquotedblright{}.) Say it once more.
\end{tcolorbox}
\section[xiōngmèi]{xiōngmèi — elder brother and younger sister}

\label{lesson:ZH-C508-xiongmei}



\begin{quote}
Before the new one: say the Mandarin for one's elders, then the Mandarin for sisters.
\end{quote}

\subsection*{You'll want to know: xiōngmèi}
\emph{xiōngmèi} — \textquotedblleft{}elder brother and younger sister\textquotedblright{}.

\begin{quote}
\emph{tāmen xiōngmèi dōu shì dàxuésheng} — The brother and sister are both at university.
\end{quote}

Its characters come later; for now it is a word for the ear.

\begin{cousinweb}
\emph{xiōngmèi} is built from \emph{xiōng} 'an elder brother' and \emph{mèi} 'a younger sister', both inherited from Middle Chinese, \emph{xjwaeng} and \emph{mwojH} (Guangyun, in Baxter's transcription).
\end{cousinweb}

\begin{grammarlens}[title={What you've built}]
One more word for saying how it was.
\end{grammarlens}

\subsection*{Your turn}
\begin{itemize}
\raggedright
  \item \emph{Say it:} \emph{xiōngmèi}
  \item \emph{Say it:} \emph{xiōngmèi}, once more
  \item \emph{Say it:} \emph{tāmen xiōngmèi dōu shì dàxuésheng} — The brother and sister are both at university.
  \item \emph{Recall:} say the Mandarin for one's elders, then the Mandarin for sisters, then say \emph{xiōngmèi} again
\end{itemize}

\begin{tcolorbox}[breakable,colback=teal!4,colframe=teal!35!black,title={Before you move on}]
What does xiōngmèi mean? (\textquotedblleft{}Elder brother and younger sister\textquotedblright{}.) Say it once more.
\end{tcolorbox}
\section[dúshēngzǐnǚ]{dúshēngzǐnǚ — only child}

\label{lesson:ZH-C508-dushengzinu}



\begin{quote}
Before the new one: say the Mandarin for sisters, then the Mandarin for elder brother and younger sister.
\end{quote}

\subsection*{You'll want to know: dúshēngzǐnǚ}
\emph{dúshēngzǐnǚ} — \textquotedblleft{}only child\textquotedblright{}.

\begin{quote}
\emph{wǒ shì dúshēngzǐnǚ xiǎoshíhou zuì xiǎng yào yí ge dìdi} — I'm an only child; when I was little, what I wanted most was a little brother.
\end{quote}

Its characters come later; for now it is a word for the ear.

Two third tones meet in this word: zǐ is said with a rising tone before nǚ, and the spelling keeps both marks.

\begin{cousinweb}
\emph{dúshēngzǐnǚ} is built from \emph{dú} 'alone', \zh{生} 'to be born' and \zh{子女} 'sons and daughters', all inherited from Middle Chinese, \emph{duwk}, \emph{srjaeng}, \emph{tsiX} and \emph{nrjoX} (Guangyun, in Baxter's transcription). For \emph{duwk} the dictionary quotes an older one that sets it against \emph{qún} 'in flocks', the way sheep go.
\end{cousinweb}

\begin{grammarlens}[title={What you've built}]
One more word for saying how it was.
\end{grammarlens}

\subsection*{Your turn}
\begin{itemize}
\raggedright
  \item \emph{Say it:} \emph{dúshēngzǐnǚ}
  \item \emph{Say it:} \emph{dúshēngzǐnǚ}, once more
  \item \emph{Say it:} \emph{wǒ shì dúshēngzǐnǚ xiǎoshíhou zuì xiǎng yào yí ge dìdi} — I'm an only child; when I was little, what I wanted most was a little brother.
  \item \emph{Recall:} say the Mandarin for sisters, then the Mandarin for elder brother and younger sister, then say \emph{dúshēngzǐnǚ} again
\end{itemize}

\begin{tcolorbox}[breakable,colback=teal!4,colframe=teal!35!black,title={Before you move on}]
What does dúshēngzǐnǚ mean? (\textquotedblleft{}Only child\textquotedblright{}.) Say it once more.
\end{tcolorbox}
\section[qīfu]{qīfu — to bully; to pick on}

\label{lesson:ZH-C508-qifu}



\begin{quote}
Before the new one: say the Mandarin for elder brother and younger sister, then the Mandarin for an only child.
\end{quote}

\subsection*{You'll want to know: qīfu}
\emph{qīfu} — \textquotedblleft{}to bully; to pick on\textquotedblright{}.

\begin{quote}
\emph{xiǎoshíhou gēge chángcháng qīfu wǒ} — When I was little, my elder brother often picked on me.
\end{quote}

Its characters come later; for now it is a word for the ear.

\begin{cousinweb}
\emph{qīfu} is built from \emph{qī} 'to deceive' and \emph{fù} 'to bear', both inherited from Middle Chinese, \emph{khi} and \emph{bjuwX} (Guangyun, in Baxter's transcription), which glosses the first 'to deceive' and the second 'to carry' and 'to repay kindness with ingratitude'; in the word the second syllable is light.
\end{cousinweb}

\begin{grammarlens}[title={What you've built}]
One more word for saying how it was.
\end{grammarlens}

\subsection*{Your turn}
\begin{itemize}
\raggedright
  \item \emph{Say it:} \emph{qīfu}
  \item \emph{Say it:} \emph{qīfu}, once more
  \item \emph{Say it:} \emph{xiǎoshíhou gēge chángcháng qīfu wǒ} — When I was little, my elder brother often picked on me.
  \item \emph{Recall:} say the Mandarin for elder brother and younger sister, then the Mandarin for an only child, then say \emph{qīfu} again
\end{itemize}

\begin{tcolorbox}[breakable,colback=teal!4,colframe=teal!35!black,title={Before you move on}]
What does qīfu mean? (\textquotedblleft{}To bully; to pick on\textquotedblright{}.) Say it once more.
\end{tcolorbox}
\section[dǒngshì]{dǒngshì — sensible; grown-up for one's age (of a child)}

\label{lesson:ZH-C508-dongshi}



\begin{quote}
Before the new one: say the Mandarin for an only child, then the Mandarin for to pick on (someone weaker).
\end{quote}

\subsection*{You'll want to know: dǒngshì}
\emph{dǒngshì} — \textquotedblleft{}sensible; grown-up for one's age (of a child)\textquotedblright{}.

\begin{quote}
\emph{jiějie bǐ wǒ dǒngshì} — My big sister is more grown-up than me.
\end{quote}

Its characters come later; for now it is a word for the ear.

\begin{cousinweb}
\emph{dǒngshì} is built from \emph{dǒng} 'to understand' and \zh{事} 'a matter'. The second is inherited from Middle Chinese \emph{dzriH} (Guangyun, in Baxter's transcription), glossed 'to employ'; the Guangyun has no entry for the first, so its history is not settled.
\end{cousinweb}

\begin{grammarlens}[title={What you've built}]
One more word for saying how it was.
\end{grammarlens}

\subsection*{Your turn}
\begin{itemize}
\raggedright
  \item \emph{Say it:} \emph{dǒngshì}
  \item \emph{Say it:} \emph{dǒngshì}, once more
  \item \emph{Say it:} \emph{jiějie bǐ wǒ dǒngshì} — My big sister is more grown-up than me.
  \item \emph{Recall:} say the Mandarin for an only child, then the Mandarin for to pick on (someone weaker), then say \emph{dǒngshì} again
\end{itemize}

\begin{tcolorbox}[breakable,colback=teal!4,colframe=teal!35!black,title={Before you move on}]
What does dǒngshì mean? (\textquotedblleft{}Sensible; grown-up for one's age (of a child)\textquotedblright{}.) Say it once more.
\end{tcolorbox}
\section[(dialogue)]{Looking back: from Songs to Brothers and Sisters}

\label{lesson:ZH-R508-review-504-508}



\begin{quote}
Say the Mandarin for the words of a song and a nursery rhyme.
\end{quote}

\subsection*{Your turn}
\begin{itemize}
\raggedright
  \item \emph{Say it:} the words of a song, a nursery rhyme, the national anthem, a folk song, to perform a song (on stage)
  \item \emph{Say it:} the art of dance, a dance party, a dance partner, a stage (in a theatre), to rehearse
  \item \emph{Say it:} a band (of musicians), to perform on an instrument, a violin, Beijing opera, a fan (of a singer)
  \item \emph{Say it:} brothers, a father and son, a father and daughter, a married couple, one's elders
  \item \emph{Say it:} sisters, elder brother and younger sister, an only child, to pick on (someone weaker), sensible (of a child)
  \item \emph{Read it:} \zh{歌} \zh{舞} \zh{父} — each one aloud, then a word that holds it
\end{itemize}

\begin{tcolorbox}[breakable,colback=teal!4,colframe=teal!35!black,title={Before you move on}]
The words of a song? (\emph{gēcí}.) A nursery rhyme? (\textbf{\zh{儿歌}}, \emph{érgē}.) The national anthem? (\textbf{\zh{国歌}}, \emph{guógē}.) A folk song? (\emph{míngē}.) To perform a song (on stage)? (\emph{yǎnchàng}.) The art of dance? (\emph{wǔdǎo}.) A dance party? (\textbf{\zh{舞会}}, \emph{wǔhuì}.) A dance partner? (\emph{wǔbàn}.) A stage (in a theatre)? (\emph{wǔtái}.) To rehearse? (\emph{páiliàn}.) A band (of musicians)? (\emph{yuèduì}.) To perform on an instrument? (\emph{yǎnzòu}.) A violin? (\emph{xiǎotíqín}.) Beijing opera? (\emph{Jīngjù}.) A fan (of a singer)? (\emph{gēmí}.) Brothers? (\emph{xiōngdì}.) A father and son? (\textbf{\zh{父子}}, \emph{fùzǐ}.) A father and daughter? (\textbf{\zh{父女}}, \emph{fùnǚ}.) A married couple? (\emph{fūqī}.) One's elders? (\emph{zhǎngbèi}.) Sisters? (\emph{jiěmèi}.) Elder brother and younger sister? (\emph{xiōngmèi}.) An only child? (\emph{dúshēngzǐnǚ}.) To pick on (someone weaker)? (\emph{qīfu}.) Sensible (of a child)? (\emph{dǒngshì}.)
\end{tcolorbox}
